\documentclass[border=5pt]{standalone}
\usepackage{tikz}
\usetikzlibrary{calc, angles, quotes}

\begin{document}
\begin{tikzpicture}
    % Define coordinates for the triangle vertices
    \coordinate (A) at (2, 5);
    \coordinate (B) at (0, 0);
    \coordinate (C) at (6, 0);
    
    % Define M as the midpoint of AB (standard geometry setup assumption)
    \coordinate (M) at ($(A)!0.5!(B)$);
    
    % Draw the original triangle and segment CM
    \draw (A) -- (B) -- (C) -- cycle;
    \draw (C) -- (M);
    
    % Label the original points
    \node[left] at (A) {$A$};
    \node[below left] at (B) {$B$};
    \node[below right] at (C) {$C$};
    \node[above left] at (M) {$M$};
    
    % --- EDIT: Perpendicular from B to CM ---
    
    % Calculate the foot of the perpendicular H from B to line CM
    \coordinate (H) at ($(C)!(B)!(M)$);
    
    % Draw segment BH
    \draw (B) -- (H);
    
    % Label H
    \node[above right] at (H) {$H$};
    
    % Draw the right angle mark at H
    \draw pic [draw, angle radius=2.5mm] {right angle = B--H--C};
    
\end{tikzpicture}
\end{document}
